% Multi-region transfer experiments (US-073, EPIC 11 Paper Track).
% \input-able section: no \documentclass / \begin{document} here.
% Prose in English (the manuscript is EN). Every number is rendered from the real
% EPIC 12 artefacts by scripts/build_us073_transfer_figures.py; nothing is hand-typed.
% Tables : paper/tables/us-073-transfer/{sen4agrinet_domain_gap,eurocropsml_kshot}.tex
% Figures: paper/figures/us-073-transfer/{kshot_curve,worldcereal_tropical_transfer}.{png,svg}
%          paper/figures/us-075/sen4agrinet_per_class_iou_f1.png
%          paper/figures/us-076/eurocropsml_per_class_f1_vs_k.png
%          figures/us-070/eurocropsml_ae_vs_s2_es.png
% Required packages in the host preamble: booktabs, graphicx, amsmath, hyperref, url.

\section{Transferencia multirregión: entrenar en Francia, extender a otros lugares}
\label{sec:multiregion}

Un copiloto de mapeo de cultivos entrenado solo con etiquetas francesas
(PASTIS-R) es, por construcción, un experto en Francia. Esta sección
convierte esa limitación en una capacidad medida: documentamos la
\emph{receta} ``entrenar en Francia $\rightarrow$ extender a otros
lugares'' a lo largo de dos rutas de transferencia reales que se
ejecutaron de extremo a extremo y se versionaron con DVC, más una
extensión tropical explícitamente fuera de alcance. La tesis que defendemos es
deliberadamente conservadora: el modelo \emph{recibe nuevas clases de
otros datasets mediante adaptación few-shot}, y la brecha de dominio
franco-ibérica es \emph{catastrófica en zero-shot}. En consecuencia, el
\textbf{$\Delta$mIoU entre zero-shot y few-shot es el entregable
científico}, no una exactitud absoluta alta. Esto es consistente con la
transferibilidad espacial limitada de los embeddings congelados de
observación de la Tierra reportada en ``Harvesting
AlphaEarth''~\cite{harvesting2026alphaearth}: hay que presupuestar para
un ajuste fino few-shot en lugar de prometer magia zero-shot.

Todos los embeddings son AlphaEarth Foundations~\cite{brown2025alphaearth}
(Google DeepMind), accedidos a través del activo de Google Earth Engine
\texttt{SATELLITE\_EMBEDDING/V1/ANNUAL}, versión de datos~1.1,
64-dimensional, CC-BY-4.0. El razonador congelado que convierte las
salidas del modelo en explicaciones es
Gemini~3.5~Flash~\cite{google2025gemini25} (el backend de nube desplegado);
siguiendo el patrón ``Be My Eyes''~\cite{huang2025bemyeyes}, nuestros
modelos entrenados son el \emph{perceptor} y el LLM es un \emph{razonador}
congelado para la comunicación y la soberanía de datos, nunca un
clasificador de píxeles.

La Figura~\ref{fig:aoi-italy} sitúa una de las áreas de interés europeas
del proyecto --- la región agrícola italiana usada para la destilación
contrastiva de FarSLIP (\texttt{farslip-clip-italy-v1}) --- dentro de la
huella multirregión que abarca Francia (PASTIS-R), Cataluña
(Sen4AgriNet), los países de EuroCropsML (Estonia/Letonia/Portugal) y las
regiones tropicales de WorldCereal (Brasil, India).

\begin{figure}[htbp]
  \centering
  \includegraphics[width=0.75\linewidth]{figures/us-070/aoi_italy_es.png}
  \caption{Área de interés italiana (Pianura Padana, valle del Po) usada
  para la destilación contrastiva de FarSLIP
  (\texttt{farslip-clip-italy-v1}), un nodo de la huella multirregión de
  AgroSatCopilot. Los píxeles reales muestreados de AlphaEarth se colorean
  según su embedding anual de 64 dimensiones sobre un mapa base de
  OpenStreetMap, para ubicar el AOI y mostrar que el embedding varía
  espacialmente en la región. AlphaEarth
  \texttt{SATELLITE\_EMBEDDING/V1/ANNUAL} v1.1
  (CC-BY-4.0)~\cite{brown2025alphaearth}; mapa base \textcopyright{}
  contribuidores de OpenStreetMap.}
  \label{fig:aoi-italy}
\end{figure}

\subsection{Armonización del espacio de etiquetas (HCAT v3)}
\label{sec:multiregion:labelspace}

La transferencia entre datasets requiere un espacio de etiquetas común,
de lo contrario los cabezales por región pelean por gradientes
incompatibles. Mapeamos las 18 clases agronómicas de PASTIS-R (el espacio
contiguo \texttt{semantic18}, ids $0\ldots17$) a los códigos de hoja de
HCAT~v3 (10 dígitos, prefijo \texttt{33}~$=$~\texttt{crop\_type}) y los
colapsamos, vía el nivel de grupo de la jerarquía HCAT, a \textbf{11
macroclases canónicas} (10 grupos de cultivo más una clase unificada
\texttt{void}/fondo), bien dentro del rango solicitado de 10--20. Se
verifica que cada código de hoja HCAT aparezca textualmente en la
referencia HCAT~v3 de EuroCrops
(\texttt{data/reference/eurocrops\_hcat3.csv}); ningún código se inventa.
El mapeo y su regla de colapso se documentan en
\texttt{docs/data/hcat\_crosswalk.md} y se materializan como
\texttt{data/reference/hcat\_crosswalk.parquet}.

Tres clases de PASTIS no tienen una hoja HCAT 1:1 y se marcan
\texttt{approx} como decisiones agronómicas explícitas (no fallos de
coincidencia): \emph{Cereal mixto} mapea al nodo de grupo L2
\texttt{cereal}; \emph{Frutas/verduras/flores} mapea a la hoja
representativa dominante \texttt{fresh\_vegetables}; \emph{Forraje
leguminoso} mapea a \texttt{legumes\_harvested\_green}. La cola larga se
mitiga con el colapso: \emph{Prado} ($\sim$45\% de las parcelas de
cultivo) se aísla como \texttt{grassland}, y los ocho cereales hermanos
(incluidos los casi inseparables trigos blando y duro de invierno) se
fusionan en \texttt{cereals}, reduciendo la varianza interclase que
deprimía el F1 macro de 18 clases. La convención entre datasets es
uniforme: \texttt{background}/\texttt{void} $\rightarrow$
\texttt{ignore\_index = 255} (excluido de la matriz de confusión),
mientras que una clase presente en un dataset pero ausente en otro se
trata como \texttt{null-class} (una etiqueta desconocida, de etiquetado
parcial), nunca como un negativo duro.

\subsection{Ruta densa: Sen4AgriNet Cataluña (Francia $\rightarrow$ España)}
\label{sec:multiregion:dense}

El mejor segmentador francés individual, TSViT~\cite{tarasiou2023tsvit}
con la rama fenológica (TSViT-pheno, fold-4 mIoU~0.6253 / F1~0.7500 en
PASTIS-R~\cite{garnot2021pastis}), es el modelo origen. Lo ajustamos
finamente sobre un pequeño subconjunto de
Sen4AgriNet~\cite{sykas2022sen4agrinet} en la tesela 31TCG de Cataluña
(Sentinel-2, años 2019/2020),
proyectando ambos cabezales sobre el espacio macro HCAT compartido. El
subconjunto es intencionalmente diminuto (10 parches de entrenamiento /
90 teselas; 20 parches de validación / 180 teselas), versionado en
DVC/GCS (\texttt{data/sen4agrinet.dvc}); nunca descargamos el dataset
completo de múltiples terabytes.

La Tabla~\ref{tab:sen4agrinet_domain_gap} reporta el resultado. La
transferencia zero-shot arroja una mIoU de exactamente~$0.0000$. Esto es
real, no un bug: el TSViT-pheno francés proyectado $18\rightarrow$macro
predice las teselas de validación de Cataluña casi exclusivamente como
\texttt{vegetables} ($\sim$567k px) y \texttt{grassland} ($\sim$416k px),
dos clases cuyo soporte es \emph{cero} en la verdad de terreno española
(dominada en cambio por \texttt{cereals}, \texttt{legumes\_fodder},
\texttt{vineyard} y \texttt{oilseed\_industrial}). Con una diagonal de
confusión toda en cero, la mIoU es exactamente~0 --- una brecha de
dominio franco-ibérica catastrófica de manual. También encontramos que
las etiquetas de Sen4AgriNet Cataluña llevan códigos de cultivo S4A fuera
de la nomenclatura PASTIS-18; estos se enrutan honestamente a
\texttt{ignore\_index = 255}, pero la verdad de terreno restante
\texttt{cereals}/\texttt{legumes}/\texttt{vineyard}/\texttt{oilseed} es
suficiente para una medición válida. Ajustar finamente el cabezal macro
sobre los 10 parches españoles recupera una mIoU few-shot de~$0.2468$
(F1-macro~$0.3005$, exactitud de píxel~$0.9179$), es decir
$\Delta$mIoU~$= +0.2468$ en 40 épocas. Este modelo ajustado es el
segmentador ``franco-ibérico'' que añade el comportamiento de olivo,
viñedo, sorgo y trigo duro al origen francés. La corrida completa se
registra en
\texttt{reports/segmentation/sen4agrinet\_transfer\_result.json} y se
reproduce con el notebook ejecutado
\texttt{notebooks/segmentation/5c\_transfer\_sen4agrinet.ipynb}.

\input{tables/us-073-transfer/sen4agrinet_domain_gap.tex}

\subsection{Transfer dentro del mismo país: Francia $\rightarrow$ Bretaña (BreizhCrops)}
\label{sec:multiregion:breizhcrops}

Antes de cruzar una frontera, planteamos la pregunta de transfer más barata posible: ¿el
modelo aguanta dentro del \emph{mismo país}? Bretaña es francesa pero agronómicamente
distinta --- clima oceánico, otro calendario de cultivo y otro año ---, así que aísla el
cambio de \emph{región} del cambio de \emph{país}. Entrenamos el clasificador tabular del
proyecto (XGBoost \texttt{multi:softprob} sobre los 185 rasgos compartidos por parcela)
en PASTIS-R (Francia, 2019) y lo evaluamos \emph{directo}, sin reentrenar, sobre
BreizhCrops~\cite{russwurm2020breizhcrops} (Bretaña, 2017; regiones
\texttt{frh01}$+$\texttt{frh04}), sobre los siete cultivos que ambos datasets comparten
($64{,}121$ parcelas de entrenamiento / $2{,}145$ de prueba). BreizhCrops no trae embedding
AlphaEarth, de modo que esta es una prueba con rasgos crudos, no una evaluación de Voting-3.

El F1-macro del transfer directo es de solo $0.2057$ (\emph{accuracy} $0.2956$): incluso
dentro de Francia, cambiar de región es un salto real. El desglose por cultivo es la señal
honesta (se genera desde \texttt{data/transfer/pastis\_to\_breizhcrops.parquet}): el transfer
\emph{aguanta} donde el cultivo se ve fenológicamente parecido entre regiones --- colza
$0.70$, pradera permanente $0.51$ --- y \emph{se hunde} en los cereales cuya dinámica
intra-anual difiere --- maíz $0.14$, cebada $0.07$, trigo blando $0.01$, girasol $0.00$.
Este resultado dentro de Francia es la ilustración más fuerte y más barata del caveat de
transferibilidad limitada (Sec.~\ref{sec:multiregion:caveat}): el fallo lo impulsa la región
y su fenología, no solo el cambio de país.

\subsection{Ruta tabular few-shot: EuroCropsML (transnacional)}
\label{sec:multiregion:fewshot}

La ruta densa mide una brecha de dominio intraeuropea a nivel de píxel;
la ruta tabular mide cuántas etiquetas locales cierran una brecha
\emph{transnacional}. Dos aclaraciones factuales son obligatorias.
Primero, \textbf{Francia no forma parte de EuroCropsML}: el
benchmark~\cite{reuss2025eurocropsml} cubre Estonia (EE), Letonia (LV) y
Portugal (PT), de modo que el protocolo real es
LV$+$PT~$\rightarrow$~EE y LV~$\rightarrow$~EE (Reuss et al.
2025, Tabla~II), no ``Francia~$\rightarrow$~Estonia''.
Segundo, EuroCropsML \textbf{no lleva embedding de AlphaEarth}: cada dato
es una serie temporal de mediana anual Sentinel-2 L1C (13 bandas, B10
excluida) por parcela. Por ello reutilizamos la \emph{receta} de nuestro
baseline tabular AlphaEarth (un XGBoost \texttt{multi:softprob} sobre un
vector por parcela de ancho fijo) pero construimos ese vector a partir de
la propia serie Sentinel-2 de la parcela; la variante
AlphaEarth-vía-GEE para estos países queda como trabajo futuro. Las
etiquetas se alinean al mismo espacio macro HCAT (el código
\texttt{EC\_hcat\_c} es un código de hoja HCAT), arrojando 8 macroclases
en este subconjunto, versionado como
\texttt{data/transfer/eurocropsml.dvc}.

La Tabla~\ref{tab:eurocropsml_kshot} y la
Figura~\ref{fig:eurocropsml_kshot} dan la curva de F1-macro-frente-a-$k$
(media~$\pm$~std sobre tres semillas) para $k \in
\{1,5,10,20,100,200,500\}$ shots por clase objetivo, en tres escenarios:
LV$+$PT~$\rightarrow$~EE con preentrenamiento de origen,
LV~$\rightarrow$~EE con preentrenamiento de origen, y una referencia sin
preentrenamiento. La lectura es la prevista para una brecha de dominio:
en $k{=}1$ el modelo sin preentrenamiento es casi aleatorio
($\sim$0.015 F1-macro) mientras que el preentrenamiento de origen ya
alcanza $\sim$0.32; a medida que $k$ crece las curvas convergen, lo que
confirma que un puñado de etiquetas locales --- no un zero-shot mágico
--- es lo que cierra la brecha. La curva es evidencia del costo de
transferencia medido, nunca una afirmación de exactitud zero-shot, y se
calcula estrictamente sobre las particiones reales de EuroCropsML.

\input{tables/us-073-transfer/eurocropsml_kshot.tex}

\begin{figure}[t]
  \centering
  \includegraphics[width=0.9\linewidth]{figures/us-073-transfer/kshot_curve_es.png}
  \caption{Curva few-shot transnacional de EuroCropsML
  (LV$[+$PT$]\rightarrow$EE). F1-macro sobre el conjunto de consulta de
  Estonia frente a $k$ shots por clase, media sobre tres semillas con
  barras de error std. El preentrenamiento de origen domina en $k$ bajo y
  todas las curvas convergen a medida que crece la supervisión local: la
  brecha la cierran los datos few-shot, no la transferencia zero-shot.
  EuroCropsML (Reuss et al. 2024, CC-BY-SA-4.0).}
  \label{fig:eurocropsml_kshot}
\end{figure}

\subsection{¿Ayuda el modelo de fundación? AlphaEarth vs.\ Sentinel-2 crudo}
\label{sec:multiregion:fm-vs-s2}

La curva anterior se construye a partir de series Sentinel-2 crudas por
parcela porque EuroCropsML no trae embedding. Una pregunta natural y
decisiva es si el modelo de fundación AlphaEarth --- el mismo embedding
anual de 64 dimensiones que sustenta el resto de este trabajo ---
realmente \emph{cierra la brecha de dominio transnacional con menos
etiquetas locales} que las bandas espectrales crudas. Para responderla
reejecutamos el protocolo few-shot idéntico LV$\rightarrow$EE (XGBoost
\texttt{multi:softprob}, tres semillas, nueve macroclases) sobre dos
espacios de características: la propia serie Sentinel-2 de la parcela, y
el embedding zonal de AlphaEarth obtenido para las mismas
$42{,}473$ parcelas de EuroCropsML a través de Google Earth
Engine~(\texttt{SATELLITE\_EMBEDDING/V1/ANNUAL}, v1.1).

La Figura~\ref{fig:eurocropsml_ae_vs_s2} reporta la comparación directa.
\textbf{AlphaEarth gana en cada $k$.} La ventaja es mayor exactamente
donde importa --- en el régimen de pocas etiquetas:
$\Delta$F1-macro~$=+0.111$ en $k{=}1$ ($0.432$ vs.\ $0.321$) y $+0.087$
en $k{=}5$ ($0.436$ vs.\ $0.349$), encogiéndose a $+0.014$--$0.035$ una
vez que hay cientos de shots locales disponibles y ambos espacios tienen
suficiente supervisión para saturarse cerca de $0.53$ en $k{=}2000$. En
otras palabras, \emph{el modelo de fundación cierra la misma brecha de
dominio con aproximadamente un orden de magnitud menos de etiquetas
locales} que las bandas crudas: un solo shot por clase sobre AlphaEarth
ya iguala lo que Sentinel-2 crudo necesita $\sim$20 shots para alcanzar.
Este es el resultado de transferencia más positivo de la sección y es la
contraparte empírica del presupuesto few-shot advertido en ``Harvesting
AlphaEarth''~\cite{harvesting2026alphaearth}: el FM no entrega magia
zero-shot, pero hace el inevitable presupuesto few-shot dramáticamente
más barato. Todas las cifras se renderizan desde
\texttt{data/transfer/eurocropsml\_alphaearth\_vs\_s2\_delta.parquet}.

\begin{figure}[t]
  \centering
  \includegraphics[width=0.9\linewidth]{figures/us-070/eurocropsml_ae_vs_s2_es.png}
  \caption{Embedding del modelo de fundación AlphaEarth vs.\ Sentinel-2
  crudo en la tarea few-shot transnacional EuroCropsML LV$\rightarrow$EE
  (F1-macro vs.\ $k$, media sobre tres semillas). AlphaEarth domina en
  cada $k$, con el mayor margen en el régimen de pocas etiquetas
  ($+0.111$ en $k{=}1$, $+0.087$ en $k{=}5$), convergiendo a medida que
  ambos espacios se saturan cerca de $0.53$ en $k{=}2000$: el modelo de
  fundación cierra la brecha de dominio con aproximadamente un orden de
  magnitud menos de etiquetas locales. Embeddings para las
  $42{,}473$ parcelas obtenidos vía Google Earth Engine
  (\texttt{SATELLITE\_EMBEDDING/V1/ANNUAL} v1.1,
  CC-BY-4.0)~\cite{brown2025alphaearth}; etiquetas EuroCropsML (Reuss et
  al.\ 2024, CC-BY-SA-4.0)~\cite{reuss2025eurocropsml}.}
  \label{fig:eurocropsml_ae_vs_s2}
\end{figure}

\subsection{Anatomía por clase: la cardinalidad domina la fragilidad}
\label{sec:multiregion:perclass}

Las puntuaciones promediadas en macro ocultan qué cultivos transfieren y
cuáles colapsan. Abrir la matriz de confusión por clase, en ambas rutas
reales, expone un único factor dominante:
\textbf{la cardinalidad por clase, no el desplazamiento de dominio
abstracto, gobierna la fragilidad de la transferencia}.
En la ruta tabular EuroCropsML LV$\rightarrow$EE
(Figura~\ref{fig:eurocropsml_perclass}, izquierda), \texttt{grassland} ya
está \emph{saturada} en F1~$\approx0.90$ desde $k{=}10$ --- la cobertura
de suelo báltica dominante no tiene brecha que cerrar --- y
\texttt{cereals} ($0.82\rightarrow0.84$) y
\texttt{oilseed\_industrial} ($0.76\rightarrow0.81$) transfieren casi
gratis. En contraste, \texttt{potato} ($0.35\rightarrow0.50$) y
\texttt{legumes\_fodder} ($0.03\rightarrow0.29$) tienen una brecha real
que se cierra solo con cientos de etiquetas locales, mientras que
\texttt{sugar\_beet} ($\sim$16 parcelas de prueba), \texttt{orchard}
($\sim$52) y \texttt{vegetables} ($\sim$66) nunca superan F1~$0.3$ ni
siquiera en $k{=}500$: su soporte diminuto deja el espacio de
características demasiado disperso para que el $k$-shot lo rellene. La
misma firma aparece en la ruta densa de Sen4AgriNet
(Figura~\ref{fig:sen4agrinet_perclass}, derecha): few-shot
\texttt{cereals} se recupera por completo (IoU~$0.923$, F1~$0.960$),
\texttt{vineyard} muestra el patrón de manual de
alta-precisión/bajo-recall (precisión~$0.955$ pero recall~$0.151$ --- su
firma es confiable pero diez parches no la cubren), y
\texttt{oilseed\_industrial} ($2067$~px) y
\texttt{potato} ($164$~px) fallan de plano (IoU~$0.000$) incluso tras el
ajuste fino. También medimos la transferencia \emph{negativa} por clase:
el preentrenamiento en la vecina Letonia ayuda a los cultivos de
fenología compartida en $k{=}10$ (\texttt{cereals} $+0.343$,
\texttt{oilseed}~$+0.302$, \texttt{grassland}~$+0.285$) pero
\emph{perjudica} a \texttt{legumes\_fodder} ($-0.144$) y
\texttt{vegetables} ($-0.064$), de modo que ``preentrenar en el vecino''
es una decisión por cultivo, no global. Las tablas completas están en
\texttt{docs/transfer/analisis-clases-transfer.md}.

\begin{figure}[t]
  \centering
  \includegraphics[width=0.49\linewidth]{figures/us-076/eurocropsml_per_class_f1_vs_k_es.png}
  \hfill
  \includegraphics[width=0.49\linewidth]{figures/us-075/sen4agrinet_per_class_iou_f1_es.png}
  \caption{Anatomía de la transferencia por clase. \emph{Izquierda:}
  EuroCropsML LV$\rightarrow$EE F1 por macroclase frente a $k$ --- los
  cultivos robustos (grassland, cereals, oilseed) se saturan con pocas o
  ninguna etiqueta local, mientras que los cultivos de baja cardinalidad
  (sugar\_beet, orchard, vegetables) se mantienen por debajo de F1~$0.3$
  independientemente de $k$. \emph{Derecha:} Sen4AgriNet
  Francia$\rightarrow$Cataluña IoU/F1 por macroclase zero-shot vs.\
  few-shot --- cereals se recupera por completo, vineyard es de
  alta-precisión/bajo-recall, y las clases minoritarias (oilseed,
  potato) se mantienen en $0$. En ambas rutas las peores clases son las
  de menor cardinalidad: el soporte, no el desplazamiento de dominio en
  abstracto, impulsa la fragilidad. EuroCropsML
  (CC-BY-SA-4.0)~\cite{reuss2025eurocropsml}; Sen4AgriNet
  (CC-BY-SA-4.0)~\cite{sykas2022sen4agrinet}.}
  \label{fig:eurocropsml_perclass}
  \label{fig:sen4agrinet_perclass}
\end{figure}

\subsection{Fuera de Europa: transferencia tropical con WorldCereal (Brasil e India)}
\label{sec:multiregion:tropical}

Probamos si la receta sobrevive a un salto a los trópicos, donde la
fenología, el calendario y el propio catálogo de cultivos difieren de
Europa. Usando la referencia ESA
WorldCereal~\cite{vantricht2023worldcereal} (activo GEE
\texttt{ESA/WorldCereal/2021/MODELS/v100}, CC-BY-4.0) obtuvimos
embeddings de AlphaEarth vía GEE para dos regiones tropicales --- el
Cerrado brasileño y Karnataka, India --- y ejecutamos dos experimentos:
una transferencia \emph{zero-shot} Europa$\rightarrow$Brasil usando el
clasificador europeo directamente, y el mismo protocolo few-shot que en
otros lugares.

El resultado zero-shot es inequívoco y negativo: la frontera de decisión
europea \textbf{no se traduce}. El maíz, la única clase compartida con la
nomenclatura PASTIS-18, colapsa a F1~$=0.0095$ en Brasil ---
efectivamente cero. Las clases tropicales restantes (soja, tierra de
cultivo genérica, no-cultivo) no tienen contraparte europea en absoluto,
de modo que no pueden predecirse por construcción; no hay mapeo que
convierta un cabezal francés en un cabezal del Cerrado. Esta es la
declaración tropical más fuerte de la salvedad de transferibilidad
limitada~\cite{harvesting2026alphaearth}.

El few-shot recupera con fuerza (Figura~\ref{fig:worldcereal_tropical}).
Con solo $k{=}20$ shots por clase, Brasil alcanza F1-macro~$=0.626$,
alrededor del $90\%$ del techo en dominio ($0.724$), y sigue subiendo a
$0.711$ en $k{=}200$; India es más difícil pero sigue la misma forma
($0.495$ en $k{=}20$, $0.656$ en $k{=}200$). La lección es consistente
con las rutas europeas: un embedding global congelado más unas pocas
docenas de etiquetas locales por clase transfiere mucho mejor que una
frontera trasplantada, pero \emph{las etiquetas locales son
obligatorias}. Crucialmente, las clases de cultivo tropicales \textbf{no}
mapean al espacio de hojas europeo PASTIS-18 --- solo el maíz se solapa
--- de modo que un copiloto verdaderamente global necesita datos few-shot
etiquetados localmente en cada zona climática, no meramente más regiones
europeas. Las cifras se renderizan desde
\texttt{data/transfer/worldcereal\_fewshot\_results.parquet} y
\texttt{worldcereal\_fewshot\_india.parquet}.

\begin{figure}[t]
  \centering
  \includegraphics[width=0.9\linewidth]{figures/us-073-transfer/worldcereal_tropical_transfer_es.png}
  \caption{Transferencia tropical con embeddings AlphaEarth de
  WorldCereal. El zero-shot Europa$\rightarrow$Brasil falla (maíz
  F1~$=0.0095$; la frontera de decisión europea no se traduce), mientras
  que el few-shot recupera el grueso del desempeño en dominio: el Cerrado
  brasileño alcanza F1-macro~$0.626$ en $k{=}20$ ($\sim$90\% del techo en
  dominio $0.724$) y $0.711$ en $k{=}200$; Karnataka, India sigue la
  misma forma ($0.495$ en $k{=}20$). Las clases de cultivo tropicales no
  mapean al espacio de hojas europeo PASTIS-18 (solo el maíz se solapa),
  de modo que la cobertura global requiere supervisión few-shot local por
  zona climática, no más regiones europeas. Embeddings AlphaEarth
  obtenidos vía Google Earth Engine
  (\texttt{SATELLITE\_EMBEDDING/V1/ANNUAL} v1.1,
  CC-BY-4.0)~\cite{brown2025alphaearth}; etiquetas de referencia ESA
  WorldCereal
  (\texttt{ESA/WorldCereal/2021/MODELS/v100}, CC-BY-4.0)~\cite{vantricht2023worldcereal}.}
  \label{fig:worldcereal_tropical}
\end{figure}

\subsection{Un único clasificador multirregión: un resultado negativo honesto}
\label{sec:multiregion:multiregionmodel}

Las rutas anteriores adaptan cada una a un objetivo. Una idea más
ambiciosa --- que probamos y reportamos aquí por completo, incluidos sus
modos de fallo --- es entrenar \emph{un} clasificador sobre el embedding
de AlphaEarth agrupando \emph{todas} las regiones a la vez (PASTIS
Francia, EuroCropsML Estonia y Letonia, Sen4AgriNet Cataluña y Francia,
WorldCereal Brasil e India). La hipótesis, debida a un miembro del
equipo, es atractiva: una clase de cultivo que es \emph{débil} en una
región por falta de muestras podría ser \emph{fuerte} en otra que abunda
en ella, de modo que agrupar debería producir una taxonomía más fina y
rica que cualquier dataset individual --- por ejemplo
\texttt{spring\_barley} donde PASTIS solo dice ``cereal''. Presentamos el
resultado honestamente: \textbf{la idea funciona parcialmente --- agranda
genuinamente la taxonomía --- pero no rescata ninguna clase por encima de
un umbral de calidad de producción, y ese hallazgo negativo es en sí
mismo la contribución.} El estudio completo está en
\texttt{docs/transfer/modelo-multiregion.md}.

\paragraph{Lo que funciona: la taxonomía realmente crece ($18\rightarrow30$ hojas).}
Armonizando cada región a un espacio HCAT compartido, el vocabulario de
hojas finas aprendible se expande de las $18$ hojas de PASTIS a
\textbf{30 hojas HCAT distintas}, \emph{sin mapeos inventados}: las $12$
hojas extra las aporta EuroCropsML y son exactamente las distinciones que
PASTIS colapsa --- \texttt{spring\_barley} vs.\ \texttt{winter\_barley},
\texttt{oats}, \texttt{rye}, \texttt{spring}/\texttt{summer\_rapeseed},
\texttt{apples}, \texttt{clover}, \texttt{alfalfa\_lucerne}. Como todas
las regiones viven en el mismo espacio AlphaEarth de 64 dimensiones, la
unión es representacionalmente válida aunque los espacios de etiquetas
difieran. El cabezal fino se entrena solo sobre las $17{,}516$ parcelas
que portan una hoja fina real (PASTIS $+$ EuroCropsML); las regiones solo
de macro (Sen4AgriNet, WorldCereal) se reservan para la evaluación
macro-colapsada y nunca inyectan una pseudoclase gruesa que canibalizaría
las hojas finas.

\paragraph{Lo que no funciona: cero clases rescatadas.}
Evitamos deliberadamente reportar un F1 macro de $N$ clases (que cae
mecánicamente a medida que $N$ crece) y en su lugar contamos cuántas
clases de hoja individuales cruzan un umbral de F1, contra el baseline
solo-PASTIS (Tabla~\ref{tab:multiregion_threshold}). El modelo
multirregión coloca $7$ hojas por encima de F1~$0.70$ y $2$ por encima de
$0.85$; solo-PASTIS con el mismo presupuesto de datos coloca $6$/$2$, y
la referencia PASTIS con datos completos coloca $10$/$2$. Definiendo una
clase \emph{rescatada} como una hoja que cruza F1~$0.85$ en multirregión
y que \emph{no} existe como clase aprendible en el espacio solo-PASTIS,
el resultado es \textbf{0 clases rescatadas en el umbral 0.85}: las dos
clases que sí cruzan ($\texttt{winter\_rapeseed\_rape}$, F1~$0.948$;
\texttt{vineyard}, F1~$0.905$) ya estaban en PASTIS. La mejor hoja
genuinamente nueva exclusiva de EuroCropsML es \texttt{apples} en
F1~$0.686$, luego \texttt{spring\_rapeseed} en $0.719$; la división
\texttt{spring}/\texttt{winter\_barley} que el modelo intenta pero no
resuelve ($0.432$/$0.393$). Añadir $12$ hojas finas no añadió $12$ clases
buenas.

\begin{table}[t]
  \centering
  \caption{Clases de hoja que cruzan un umbral de F1: el modelo
  multirregión agrupado vs.\ el baseline solo-PASTIS. La taxonomía
  multirregión es más amplia (30 vs.\ 18 hojas) y no degrada la lectura
  gruesa, pero no rescata \emph{ninguna} clase por encima del umbral de
  producción de 0.85 --- las dos hojas que coloca allí ya eran aprendibles
  en PASTIS. Cifras de
  \texttt{data/transfer/multiregion\_summary.json}.}
  \label{tab:multiregion_threshold}
  \begin{tabular}{lcccc}
    \toprule
    Modelo & Hojas totales & F1$\ge0.70$ & F1$\ge0.80$ & F1$\ge0.85$ \\
    \midrule
    Multirregión (agrupado)            & 30 & \textbf{7} & 2 & 2 \\
    Solo-PASTIS (acotado, mismo presupuesto) & 18 & 6 & 3 & 2 \\
    Solo-PASTIS (completo, referencia)     & 18 & 10 & 4 & 2 \\
    \midrule
    \multicolumn{4}{l}{\emph{Clases rescatadas ($\ge0.85$, exclusivas de EuroCropsML)}} & \textbf{0} \\
    \bottomrule
  \end{tabular}
\end{table}

\paragraph{El verdadero cuello de botella: separabilidad fenológica, no regiones.}
Es importante señalar que el nivel grueso \emph{no} se degrada:
colapsando cada predicción fina a su macro HCAT, el modelo multirregión
puntúa F1-macro~$=0.658$ sobre las diez macroclases europeas, a la par
con el baseline tabular monorregión del proyecto ($0.6535$), a pesar de
predecir $30$ hojas en lugar de $18$. (Incluyendo las dos macros
tropicales que no puede predecir --- \texttt{non\_crop} y
\texttt{other\_cropland}, ambas F1~$0.0$ por falta de una hoja que
colapse a ellas --- el F1 macro de todas las clases es $0.548$.) El
diagnóstico es el hallazgo central: \textbf{el cuello de botella es la
separabilidad fenológica, no el número de regiones.} La mayoría de las
nuevas hojas de EuroCropsML (avena, centeno, trébol, alfalfa, trigo de
primavera, verduras finas) se sitúan en F1~$0.06$--$0.44$ porque son
\emph{intrínsecamente difíciles de separar en un embedding anual}:
distinguir la cebada de primavera de la cebada de invierno, o la avena
del trigo, necesita la fenología intraanual que un único vector anual de
64 dimensiones promedia y diluye. Agrupar más \emph{regiones} no ayuda a
una clase que ninguna región puede separar a esta resolución temporal; el
siguiente paso indicado son características temporales intraanuales
multifecha sobre las hojas de cereal/leguminosa, no más regiones. El
modelo sigue siendo útil como un vocabulario \emph{macro} más rico, pero
aún no como un clasificador fino de producción. Lo reportamos para que el
límite quede registrado, no oculto.

\subsection{Salvedad: transferibilidad espacial limitada}
\label{sec:multiregion:caveat}

Cada ruta real corrobora la salvedad central de ``Harvesting
AlphaEarth''~\cite{harvesting2026alphaearth}: los embeddings globales
congelados transfieren mal entre regiones agronómicas distantes sin
adaptación local. La mIoU zero-shot densa
Francia~$\rightarrow$~Cataluña de~$0.0000$ y el F1 zero-shot tropical del
maíz Europa~$\rightarrow$~Brasil de~$0.0095$ son las declaraciones
empíricas más fuertes posibles de esta brecha, mientras que las curvas
few-shot de EuroCropsML y WorldCereal cuantifican el presupuesto
necesario para salvarla --- y la comparación AlphaEarth-vs-Sentinel-2
muestra que el modelo de fundación reduce ese presupuesto en
aproximadamente un orden de magnitud en etiquetas locales ($+0.111$ de F1
en $k{=}1$). El contrapeso honesto es el clasificador multirregión
agrupado: amplía la taxonomía ($18\rightarrow30$ hojas) sin degradar la
lectura gruesa, pero rescata $0$ clases por encima de F1~$0.85$, porque
el verdadero cuello de botella es la separabilidad fenológica en un
embedding anual, no el número de regiones. Por tanto \emph{no} afirmamos
que el zero-shot funcione fuera de Francia, que agrupar regiones rescate
nuevas clases finas, que la capa LLM ``supere'' a Gemini en
clasificación (su rol es la comunicación, la explicación y la soberanía
de datos on-premises vía Qwen3-30B-A3B~\cite{yang2025qwen3} servido con
vLLM GPTQ-Int4 en una sola GPU), ni ningún F1 $\geq$0.80 en los
trópicos. Los entregables de esta sección son el $\Delta$mIoU
\emph{medido}, las curvas few-shot, la ventaja FM-vs-crudo y el hallazgo
negativo del modelo agrupado --- evidencia honesta de cómo, y hasta qué
punto, el copiloto se convierte en un experto multirregión.

% Share-alike note: figures and reports that combine derivatives of Sen4AgriNet,
% EuroCropsML and WorldCereal inherit the CC-BY(-SA)-4.0 clauses and must be redistributed
% under a compatible license (see docs/licenses/DATA_LICENSE.md). WorldCereal is CC-BY-4.0;
% Sen4AgriNet and EuroCropsML are CC-BY-SA-4.0.
